\section{Purpose, scope, and conventions}
\label{sec:introduction}

The original work motivating this project is Levent Alp\"oge's announced
proof of the Riemann Hypothesis and the Generalized Riemann Hypothesis,
developed with Claude and accompanied by a released formal
source~\cite{Alpoge,AlpogeSource}. The project brief identifies the model as
Fable~5.2; the announcement itself names Claude without specifying a
model version. We attribute the announced proof to Alp\"oge and distinguish
it from the expository reconstructions and checks carried out here.

The assertion under study is GRH for Dirichlet $L$-functions. For a
Dirichlet character $\chi$, it states that
\begin{equation}\label{eq:grh-target}
 L(s,\chi)=0,\qquad 0<\Re s<1
 \quad\Longrightarrow\quad \Re s=\frac12.
\end{equation}
The interval $0<\Re s<1$ is the open critical strip. Definitions of the
characters and the functions are given in
Section~\ref{sec:primitive-reduction}. The principal character of conductor
one gives the Riemann zeta function. No assertion about all automorphic
$L$-functions is made here.

This manuscript does not yet prove~\eqref{eq:grh-target}. It proves
reductions whose missing arithmetic inputs are stated separately.
In particular, the functional equation reduces the critical-strip
assertion to nonvanishing for primitive characters to the right of
one half. We give two complementary ways to obtain nonvanishing from
estimates: a Mellin transform of smoothed M\"obius sums, and a positivity
argument for suitable weighted prime sums. These are implications;
the required cancellation and upper bounds remain substantial mathematical
tasks.

The quasi-RH sources~\cite{OpenAISeptember,OpenAIOctober} provide useful
comparisons. The September source states a zero-free half-plane
$\Re s>7/8$, and the October source isolates a simpler route to
$\Re s>11/12$. We prove the elementary consequences of the October
second-moment input and identify exactly what its exponent permits.
Neither stated half-plane is the GRH half-plane $\Re s>1/2$.
The automorphic estimates in those sources are not treated as having
been independently proved merely because their consequences are proved
here.

For positive real $x$, complex powers always mean
$x^s=\exp(s\log x)$ with the real logarithm. The notation
$F\ll_{\mathcal P}G$, where $G\geq0$, means that $|F|\leq CG$ for a
constant $C$ depending only on the listed fixed parameters $\mathcal P$.
Every parameter declared fixed is chosen before a limit such as
$X\to\infty$. Dependence on a fixed character is allowed unless a
uniform assertion over a character family is expressly made.

We use elementary complex analysis, finite Fourier orthogonality,
unique factorization of integers and ideals, and the ordinary integration
theorems for absolutely convergent sums. The analytic normalizations
needed below are included explicitly. The mathematical exposition is
organized by logical dependence. Source-file paths, build records, and
the continuing audit are maintained outside the paper.
